\documentclass[11pt]{article}
\usepackage[margin=1in]{geometry}
\usepackage{amsmath,amssymb}
\usepackage[hidelinks]{hyperref}

\newcommand{\Buo}{\mathsf{Buo}}
\newcommand{\uo}{\mathsf{uo}}
\newcommand{\ellp}{\ell_p}
\newcommand{\ellone}{\ell_1}
\newcommand{\sgn}{\operatorname{sgn}}

\begin{document}

\section*{Human Review: Bourgain-Root Counterexamples in \texorpdfstring{$\ell_p$}{ell-p}}

\noindent Original automatic proof: \texttt{solution\_packet.pdf}

\noindent Checked date: 17 June 2026

\noindent Status: Manually checked.

\noindent Verdict: The proof appears correct.

\section*{Human Comments}

The proof appears to correctly give a negative answer to the \(\ell_p\) bullet of Question 5.2 in Kania--Swaczyna's paper, for every \(1<p<\infty\), assuming Bourgain's original \(\ell_1\) example is being used exactly as stated.

The main idea is not to use a linear embedding or a linear isomorphism between \(\ell_1\) and \(\ell_p\). Rather, the proof uses the coordinatewise nonlinear maps
\[
  \Phi_p(t)=\sgn(t)|t|^{1/p},
  \qquad
  \Psi_p(t)=\sgn(t)|t|^p.
\]
Applied coordinatewise, these maps give an order-preserving bijective correspondence between the underlying coordinatewise ordered sets of \(\ell_1\) and \(\ell_p\). This correspondence is emphatically not additive, but additivity is not what is needed here. The relevant structure is coordinatewise order, domination, and order convergence.

Starting with Bourgain's \(\ell_1\) sequence \((z_n)\), define
\[
  x_n=\Phi_p(z_n).
\]
Then \(x_n\in\ell_p\), since
\[
  \|x_n\|_p^p=\sum_j |x_n(j)|^p=\sum_j |z_n(j)|=\|z_n\|_1.
\]

The key technical point is that the root map transports order-null difference sequences from \(\ell_1\) to \(\ell_p\). The scalar estimate
\[
  |\Phi_p(s)-\Phi_p(t)|^p\leq 2^{p-1}|s-t|
\]
is enough for this. If \(a_m-b_m\to0\) in order in \(\ell_1\), then the differences are eventually dominated by positive \(\ell_1\)-vectors \(w_\alpha\downarrow0\). The estimate gives domination in \(\ell_p\) by \(2^{1-1/p}w_\alpha^{1/p}\), and these vectors decrease to \(0\) in \(\ell_p\). Hence
\[
  \Phi_p(a_m)-\Phi_p(b_m)\to0
\]
in order in \(\ell_p\).

Therefore Bourgain's subsequence-difference Cauchy condition is transported to \(\ell_p\): for every increasing subsequence \((n_k)\),
\[
  z_{n_{k+1}}-z_{n_k}\to0 \text{ in order in } \ell_1
\]
implies
\[
  x_{n_{k+1}}-x_{n_k}\to0 \text{ in order in } \ell_p.
\]
Thus \((x_n)\) is \(\Buo\)-Cauchy in \(\ell_p\).

The non-convergence argument also appears correct. If \(x_n\) were order convergent in \(\ell_p\), then it would be coordinatewise convergent and eventually dominated by some \(y\in(\ell_p)_+\). Applying the inverse coordinatewise power map gives
\[
  z_n=\Psi_p(x_n),
\]
with coordinatewise convergence and domination by \(y^p\in\ell_1\), up to finitely many initial terms. Hence \((z_n)\) would order converge in \(\ell_1\), contradicting Bourgain's example.

\section*{Review Outcome}

Archive this packet as an apparently correct counterexample, pending any later literature or author feedback. Conceptually, the proof transfers Bourgain's \(\ell_1\) obstruction to every \(\ell_p\), \(1<p<\infty\), through the nonlinear coordinatewise order isomorphism \(t\mapsto \sgn(t)|t|^{1/p}\). The map is not additive, but the argument only needs preservation of order-null dominated differences and the inverse domination argument.

\end{document}
